\chapter{学科教学知识}\label{s:pck}

每一位老师都需要三样东西：

\begin{description}

\item[\gref{g:content-knowledge}{内容知识}]
  比如怎么编程；

\item[\gref{g:general-pedagogical-knowledge}{一般教学法知识}]
  比如对学习心理学的理解；
  以及

\item[\gref{g:pedagogical-content-knowledge}{学科教学知识}]
  关于如何把某个具体概念教给某一特定受众的领域专门知识。
  在计算领域，
  学科教学知识包括教「参数如何传给函数」时该用什么例子，
  或者关于 HTML 标签嵌套最常见的误解是什么之类的东西。

\end{description}

我们可以再往这个组合里加上技术知识~\cite{Koeh2013}，
但这不改变关键的一点：
光知道这门学科、光知道怎么教书还不够——你还得知道
怎么教\emph{这门特定的}学科~\cite{Maye2004}。
所以本章会总结计算机教育研究的一些成果，
来充实你的学科教学知识储备。

和对待所有研究一样，
解读结果时需要一些谨慎：

\begin{description}

\item[随着数据增多，理论会变化。]
  计算机教育研究是一门年轻的学科：\index{计算机教育研究}
  美国工程教育学会成立于 1893 年，
  全国数学教师委员会成立于 1920 年，
  但计算机科学教师协会直到 2005 年才出现。
  尽管像 \hreffoot{http://sigcse.org/}{SIGCSE}、
  \hreffoot{http://iticse.acm.org/}{ITiCSE}
  和 \hreffoot{https://icer.hosting.acm.org}{ICER} 这样的会议不断报告新的洞见，
  我们对「学习编程」的了解，
  就是不如我们对「学习阅读、
  做一项运动、
  或做基础算术」的了解那么多。

\item[这些研究的被试大多数是 WEIRD 的：]\index{WEIRD}
  他们来自西方（Western）、受过良好教育（Educated）、工业化（Industrialized）、富裕（Rich）、民主（Democratic）的社会~\cite{Henr2010}。
  而且，
  他们大多还很年轻、身处正规课堂，
  因为这是大多数研究者最容易接触到的群体。
  我们对成年人、
  边缘化群体的成员、
  散养环境中的学习者、
  以及\grefdex{g:end-user-programmer}{终端用户程序员}{终端用户程序员}的了解要少得多，
  尽管这些人的数量要多得多。

\end{description}

如果这是一篇学术论文，
那么大多数论断前面我都会加上「某些研究似乎表明{\ldots}」之类的限定词。
但既然真实课堂里的真实教师必须做决定，
不管研究有没有明确答案，
本章给出的就是可以行动的最佳猜测，而不是面面俱到的也许。

\begin{aside}{行话}
  和任何专业一样，
  计算机教育研究也有行话。
  \gref{g:cs1}{CS1}指一门持续一学期的入门课程，
  学习者在其中第一次接触变量、循环和函数；
  而\gref{g:cs2}{CS2}指第二门课程，
  内容涵盖栈和队列这样的基础数据结构；
  \gref{g:cs0}{CS0}则指一门面向没有任何先修经验、
  也不打算马上继续学计算的人的入门课程。
  这些术语的完整定义可以在
  \hreffoot{https://www.acm.org/education/curricula-recommendations}{ACM 课程指南}里找到。
\end{aside}

\seclbl{我们现在在教他们什么？}{s:pck-now}

关于编程集训营和其他散养式项目教了什么，人们知之甚少，
部分原因是很多机构不愿意分享自己的课程。
我们对院校在教什么了解得更多一些~\cite{Luxt2017}：

\begin{longtable}{llll}
\textbf{主题}&	\textbf{\%}&	\textbf{主题}&	\textbf{\%}	\\
编程过程	&	87\%&	数据类型	&	23\%	\\
抽象编程思维&	63\%&	输入/输出	&	17\%	\\
数据结构	&	40\%&	库	&	15\%	\\
面向对象概念	&	36\%&	变量与赋值	&	14\%	\\
控制结构	&	33\%&	递归	&	10\%	\\
运算与函数	&	26\%&	指针与内存管理&	5\%
\end{longtable}

像这样高层次的主题标签，可能掩盖许多问题。
一个更实在的结果来自\cite{Rich2017}，
它回顾了一百篇文章，
找出了小学和中学计算课的学习轨迹。
他们关于顺序、重复和条件的结果，本质上是一些集体概念图，
把许多不同教育者显性和隐性的思考汇集、梳理到一起
（\figref{f:pck-trajectory}）。

\newpage
\figpdfhere{figures/conditionals.pdf}{条件的学习轨迹（来自~\cite{Rich2017}）}{f:pck-trajectory}{}

\seclbl{他们学到了多少？}{s:pck-learning}

老师教的东西和学习者学到多少之间，可能有着天壤之别。
要探究后者，
我们就得用别的衡量方式，或者直接做研究。
先说前一种方式，
大约三分之二的高等教育学生能通过他们第一门计算课程，
具体比例随班级规模等有所不同，
但无论随时间推移、还是按语言划分，都没有显著差异~\cite{Benn2007a,Wats2014}。

先备经验对结果有什么影响？\index{先备经验的影响}
为了弄清这一点，
\cite{Wilc2018}比较了 CS1 和 CS2（见下文）中有编程先备经验和没有先备经验的新手。
他们发现，有先备经验的新手在 CS1 中比没有的多出 10\% 的分数，
但到 CS2 结束时这些差异就消失了。
他们还发现，有先备经验的女性在所有方面都胜过男性同伴，
却一直对自己的能力更缺乏信心（\secref{s:motivation-inclusivity}）。

至于对新手究竟学到多少的直接研究，
\cite{McCr2001}呈现了一项多站点国际研究，
后来被~\cite{Utti2013}重复验证。
根据第一项研究，
「{\ldots}这些令人失望的结果表明，
许多学生在入门课程结束时仍然不知道如何编程。」
更具体地说，
「对来自四所大学的 216 名学生合并样本，
按为本研究制定的一般评价标准，平均得分为 110 分中的 22.89 分。」
这个结果既说明了学生的能力，也说明了老师的期望，
但不管怎样，
我们的第一条建议是\recommendation{测量并追踪结果}，
方式要能随时间比较，
好让你知道自己的课是变得更有效还是更无效。

\seclbl{新手有哪些误解？}{s:pck-misunderstand}

\chapref{s:models}解释了为什么清除新手的误解
和教他们解决问题的策略同样重要。
新手最大的误解——在编程里有时被叫做「超级 bug」——是\index{超级 bug}
以为计算机能像另一个人那样理解人的意思~\cite{Pea1986}。
因此我们的第二条建议是\recommendation{教新手明白计算机并不理解程序},
也就是把变量命名为「cost」并不能保证它的值真的是一个成本。

\cite{Sorv2018}列出了另外四十多种老师也可以设法清除的误解，
其中许多在~\cite{Qian2017}的综述里也有讨论。
其中之一，是以为程序里的变量和电子表格里的变量工作方式一样，
也就是以为在执行完：

\begin{minted}{text}
grade = 65
total = grade + 10
grade = 80
print(total)
\end{minted}

\noindent
之后，\texttt{total} 的值会是 90 而不是 75~\cite{Kohn2017}。
这正说明新手会通过类比构建出看似合理却错误的心智模型；
其他误解还包括：

\begin{itemize}

\item
  一个变量保存着它曾被赋过的所有值的历史，
  也就是它记得自己以前的值。

\item
  两个对 \texttt{name} 或 \texttt{id} 属性有相同值的对象，
  一定是同一个对象。

\item
  函数在被定义时就被执行，
  或者按它们被定义的顺序执行。

\item
  \texttt{while} 循环的条件被不断求值，
  一旦变为假，循环就停止。
  反过来，
  \texttt{if} 语句里的条件也被不断求值，
  一旦条件变为真，其语句就被执行，
  而不管当时控制流在哪里。

\item
  赋值会搬走值，
  也就是执行 \texttt{a\ =\ b} 之后，变量 \texttt{b} 就空了。

\end{itemize}

\seclbl{新手会犯哪些错误？}{s:pck-mistakes}

新手犯的错误能告诉我们教学时该优先讲什么，
但事实证明，大多数老师并不知道不同类型错误实际有多常见。
这方面规模最大的研究是\cite{Brow2017}，
它发现引号和括号不匹配是新手 Java 程序里最常见的一类错误，
但也是最容易修复的；
而有些错误（比如把 \texttt{if} 的条件放在 \texttt{\{\}} 里而不是 \texttt{()} 里）
通常只犯一次。
不出所料，
会产生编译器错误的错误，比不产生的要修得快得多。
但有些错误，
会被犯很多次，
比如用错误的参数调用方法
（例如传了一个字符串而不是整数）。

\begin{aside}{不对与没做完}
  这类研究中的一个难点，是区分错误和未完成的工作。
  例如，
  一个空的 \texttt{if} 语句，或一个定义了但还没用过的方法，
  可能是代码尚未写完的迹象，而不是错误。
\end{aside}

\cite{Brow2017}还把新手实际犯的错误
和他们的老师以为他们会犯的错误做了比较。
他们发现，
「{\ldots}教育者对哪些错误最常出现只形成了很弱的共识，
他们的排名与{\ldots}数据里的学生只有中等程度的对应，
而且教育者的经验对这种一致程度没有影响。」
例如，
把 \texttt{=}（赋值）和 \texttt{==}（相等）搞混，
并没有大多数老师以为的那么常见。

\begin{aside}{不只是在代码里}
  \cite{Park2015}在一门入门网页开发课程中，从一个在线 HTML 编辑器收集了数据，
  发现有将近所有学习者都犯过语法错误，而且过了好几周仍未解决。
  这些错误里有 20\% 与相对复杂的规则有关，
  这些规则规定 HTML 元素\emph{何时}可以嵌套在彼此之中；
  但 35\% 与更简单的、决定 HTML 元素\emph{如何}嵌套的标签语法有关。
  许多老师倾向于说
  「可规则很简单啊」，
  这正好是\chapref{s:memory}里讨论过的专家盲区的一个例子{\ldots}
\end{aside}

\seclbl{新手是怎么编程的？}{s:pck-programming}

\cite{Solo1984,Solo1986}开创了对新手和专家编程策略的探索。
关键发现是：专家既知道「什么」，也知道「怎么」，
也就是他们既明白程序里该放什么，
\emph{又}有一套程序模式或计划来指导自己构建程序。\index{程序模式}
新手两者都缺，
但大多数老师只关注前者，
尽管 bug 往往是因为没有解决问题的策略，
而不是因为对语言的知识不够。
近期研究显示了按特定顺序教四项不同技能的有效性~\cite{Xie2019}：

\begin{longtable}{lll}
&	\textbf{代码语义}&	\textbf{与目标相关的模板}	\\
\textbf{读}&	1.\ 读代码并预测行为&	3.\ 认出模板及其用途	\\
\textbf{写}&	2.\ 写出正确的语法&	4.\ 用模板达成目标
\end{longtable}

所以我们接下来的建议是：
\recommendation{先让学习者读代码，再改代码，然后写代码}，
以及\recommendation{明确引入常见模式，并让学习者练习使用它们}。
\cite{Mull2007b}只是众多证明明确教授常见模式大有裨益的研究之一，
而把问题分解成模式，也为创建和标注子目标创造了自然的机会~\cite{Marg2012,Marg2016}。\index{标注子目标}

\seclbl{新手是怎么调试的？}{s:pck-debug}

十年前，
\cite{McCa2008}写道：
「令人惊讶的是，在大多数入门编程教科书里，
分给 bug 和调试的篇幅少得可怜。」
此后几乎没什么变化：
有几百本关于编译器和操作系统的书，
但关于调试的只有寥寥几本，\index{调试}
而我从未见过一门专门讲这个主题的本科课程。

\cite{List2004,List2009}发现，许多新手很难预测短代码片段的输出，
也很难在被告知代码应该做什么的情况下，
从一组可能性中选出正确的补全。
更晚近一些，
\cite{Harr2018}发现，能够追踪代码和能够写代码之间的差距，到 CS2 时大体上已经弥合，
但仍有差距（无论哪个方向）的新手，很可能表现不佳。

因此我们的第五条建议是\recommendation{明确教新手如何调试}。
\cite{Fitz2008,Murp2008}发现，好的调试者就是好的程序员，
但并非所有好程序员都擅长调试。
擅长调试的人会用符号调试器单步执行程序、
手工追踪执行过程、
写测试、
并频繁重读规格说明，
这些都是可以教会的习惯。
然而，
逐步追踪执行有时会被无效地使用：
例如，
新手可能在 \texttt{if}-\texttt{else} 的两个分支里放同一条 \texttt{print} 语句。
新手在试图隔离问题时，还会把其实正确的行注释掉；
老师可以故意犯这两种错误，
把它们指出来、
并加以纠正，帮新手越过它们。

教新手如何调试，也能让课堂更好管理。
\cite{Alqa2017}发现，经验更多的学习者解决调试问题的速度快得多，
但时间差异很大：
个别练习的典型范围是 4—10 分钟，
这意味着有些学习者需要比其他人多 2—3 倍的时间才能做完同样的练习。
把学得快的人的做法教给学得慢的人，
有助于让全组的整体进度更整齐。

调试依赖于能读代码，
多项研究都表明这是找 bug 最有效的一种方式~\cite{Basi1987,Keme2009,Bacc2013}。
\cite{Steg2014,Steg2016a}中开发的代码质量评分标准，
是一份很好的检查清单，
不过最好分块呈现，而不是一次全给。

让学习者读代码并总结其行为是很好的练习（\secref{s:individual-strategies}），
但往往太费时间，在课堂上不实用。
反过来，让他们在程序运行前先预测输出，
有助于强化学习（\secref{s:classroom-practices}），
也给了他们一个自然的时机去问「要是……会怎样」的问题。
老师和学习者也可以在过程中追踪变量的变化，
\cite{Cunn2017}发现这很有效（\secref{s:exercises-tracing}）。

\seclbl{那测试呢？}{s:pck-testing}
\index{测试（软件）}

新手程序员似乎和专业程序员一样不愿测试软件。
这么做有价值是毫无疑问的——\cite{Cart2017}发现，
表现好的新手花大量时间测试，
而表现差的人则花更多时间在带错误的代码上——许多老师
也要求学习者给作业写测试。
但他们做得怎么样呢？
一个答案来自~\cite{Bria2015}，
它按学习者的程序通过了多少由老师提供的测试用例来给程序打分，
反过来又按学习者写的测试用例抓住了多少蓄意埋下的 bug 来给用例打分。
他们发现，新手的测试往往覆盖率低（也就是没有测到大部分代码），
而且往往一次测很多东西，这让人很难定位错误的成因。

另一个答案来自~\cite{Edwa2014b}，
它把所有新手提交的代码中所有的 bug 合起来看，
找出哪些被新手的测试套件检测到了。
他们发现，新手的测试平均只检测出整个程序总体中存在的 13.6\% 的故障。
而且，
90\% 的新手测试都非常相似，
这说明新手写测试主要是为了确认代码在做它该做的事，
而不是为了找出它没在该做的事的情形。

教授更好测试实践的一种做法，
是通过提供一组要通过的测试来定义编程问题，
而不是通过文字描述（\secref{s:exercises-classics}）。
不过，在这么做之前，
先花点时间看看你最近为自己的代码写了多少测试，
然后决定：你是在教你认为人们应该做的事，
还是他们（和你）实际做的事。

\seclbl{语言重要吗？}{s:pck-language}
\index{积木式编程}

简短的回答是「重要」：
新手使用像 Scratch 这样的积木式工具时，学编程更快、学到更多
（\figref{f:pck-scratch}）~\cite{Wein2017}。\index{Scratch}
一个原因是，积木式系统通过消除语法错误的可能性降低了认知负荷。
另一个原因是，积木式界面鼓励探索，方式与文字不同：
和所有好工具一样，
Scratch 是可以「不小心」学会的~\cite{Malo2010}。

但在积木\emph{之后}会发生什么？
\cite{Chen2018}发现，第一门编程语言是图形化的学习者，
在入门编程课程里的成绩，比第一门语言是文本语言的学习者更高——
前提是这些语言是在青春期早期或更早引入的。
因此我们的第六条建议是：
\recommendation{让儿童和青少年从积木式界面开始}，
然后再转向文本系统。
年龄限定是有原因的，因为 Scratch 刻意做成像是给较年轻用户用的，
要说服成年人认真对待它，仍然很难。

\figimg{figures/scratch-program.png}{Scratch}{f:pck-scratch}{}

Scratch 大概比其他任何编程工具都被研究得更多。
一个例子是~\cite{Aiva2016}，
它分析了超过 25 万个 Scratch 项目，
发现（除其他之外）大约 28\% 的项目里有一些从未被调用或触发过的积木。
正如前面关于「未完成 vs 不正确」的 Java 程序的那个旁注所说，
作者推测用户可能把这些积木当作草稿本，
用来记下他们（暂时）还不想丢掉的代码片段。
另一个例子是~\cite{Grov2017,Mlad2017}，
它研究了新手在 Scratch、Logo 和 Python 中学习循环的情况。\index{Scratch}\index{Python}
他们发现，使用积木式语言而不是文本式语言时，
关于循环的误解会被降到最少。
而且，
当任务变得更复杂（比如使用嵌套循环）时，差异会变得更大。

\begin{aside}{本可不必这么难}
  编程语言的创造者没有做基本的可用性测试，从而把那些语言弄得更难学。
  例如，
  \cite{Stef2013}发现，
  「{\ldots}计算机科学里用于循环的三个最常见单词——
  \texttt{for}、\texttt{while} 和 \texttt{foreach}——
  被非程序员评为最不直观的三个选择。」
  他们的工作表明，C 风格语法（Java 和 Perl 所用的那种）
  对新手来说，和随机设计的语法一样难学，
  但 Python 和 Ruby 这类语言的语法\index{Python}\index{Ruby}
  明显更容易学，
  而那种「在加入语言之前先测试特性」的语言，语法还要更容易。
  \cite{Stef2017}是一篇有用的简短综述，讲我们关于设计编程语言实际知道什么、
  以及为什么我们相信它是对的；
  而\cite{Guzd2016}列出了面向学习者的编程语言应遵循的五条原则。
\end{aside}

\subsection*{面向对象编程}
\index{面向对象编程}

对象和类是有经验程序员的强力工具，
许多教育者主张在教编程时采用\gref{g:objects-first}{对象优先}的方法
（不过他们有时对「对象优先」到底指什么意见不一~\cite{Benn2007b}）。
\cite{Sorv2014}描述并论证了这种方法，
而\cite{Koll2015}描述了三代旨在支持新手在面向对象环境中编程的工具。

过早引入对象有一些挑战。
\cite{Mill2016b}发现，大多数使用 Python 的新手\index{Python}
都很难理解 \texttt{self}
（它指当前对象）：
他们在方法定义里漏掉它、
在引用对象属性时没有用它、
或者两样都犯。
\cite{Rago2017}在高中生身上发现了类似情况，
还发现高中老师往往对这个概念也不清楚。
权衡之后，
我们建议老师\recommendation{先教函数，而不是对象}，
也就是在学习者理解基本控制结构和数据类型之前，
不要教他们如何定义类。

\subsection*{类型声明}
\index{类型声明}

程序员们争论了几十年：变量的数据类型是否必须声明。
争论通常基于他们作为专业人士的个人经验，
而不是任何数据。
\cite{Endr2014,Fisc2015}发现，要求新手声明变量类型确实会给程序增加一些复杂度，
但它很快就能回本，因为它为该方法的用法充当了文档——尤其是，
它预先回答了「有哪些可用的东西、该怎么用」这类问题。

\subsection*{变量命名}
\index{变量命名}

\cite{Kern1999}写道：
「程序员常被鼓励无论语境如何都使用长变量名。
这是个错误：清晰往往靠简洁来实现。」
很多程序员相信这一点，
但\cite{Hofm2017}发现，在变量名里使用完整单词，
相比字母和缩写，平均让理解速度提高了 19\%。
相反，
\cite{Beni2017}发现，使用单字母变量名并不影响新手修改代码的能力。
这可能是因为他们的程序比专业人士的短，
也可能是因为某些单字母变量名带有隐含的类型和含义。
例如，
大多数程序员假定 \texttt{i}、\texttt{j} 和 \texttt{n} 是整数，
\texttt{s} 是字符串，
而 \texttt{x}、\texttt{y}、\texttt{z} 或多或少是同等地既可能是浮点数也可能是整数。

这有多重要？
\cite{Bink2012}报告说，读和理解代码，与读散文有本质的不同：
「{\ldots}源码更正式的结构和语法
让程序员能相当迅速地吸收和理解代码的各个部分，且不受风格影响。
尤其是{\ldots}信标和程序计划在理解中扮演了重要角色。」
它还发现，有经验的开发者相对不受标识符风格的影响，
所以我们的建议只是：在所有例子里使用一致的风格。
既然大多数语言都有风格指南
（例如 Python 的\hreffoot{https://www.python.org/dev/peps/pep-0008/}{PEP 8}）
以及检查代码是否遵循这些指南的工具，
我们的完整建议就是：
\recommendation{用工具确保所有代码示例遵循一致的风格}。

\seclbl{更好的错误信息有帮助吗？}{s:pck-error}
\index{错误信息}

让人看不懂的错误信息，是新手（以及有经验的程序员）挫败感的一大来源。
因此有些研究者探究了更好的错误信息是否能缓解这个问题。
例如，
\cite{Beck2016}重写了 Java 编译器的部分信息，
使得学习者看到的不是：

\begin{minted}{text}
C:\stj\Hello.java:2: error: cannot find symbol
        public static void main(string[ ] args)
^
1 error
Process terminated ... there were problems.
\end{minted}

\noindent
而是：

\begin{minted}{text}
Looks like a problem on line number 2.
If "string" refers to a datatype, capitalize the 's'!
\end{minted}

\noindent
果不其然，
拿到这些信息的新手，重复犯的错误更少，总体错误也更少。

\cite{Bari2017}更进一步，用眼动追踪表明：
尽管编译器的作者们颇有怨言，
人们确实会读错误信息——事实上，他们在这上面花掉了 13—25\% 的时间。
然而，
读错误信息被证明和读源码一样难，
而错误信息读起来有多难，能强烈预测任务表现。
因此老师应该
\recommendation{教学习者如何阅读和解读错误信息}。
\cite{Marc2011}有一份针对错误信息作答的评分标准，给这类练习打分时很有用。

\subsection*{可视化有帮助吗？}
\index{程序可视化}

把程序执行可视化，是一个经久不衰的流行想法，
像 Online Python Tutor~\cite{Guo2013}
和 \hreffoot{http://latentflip.com/loupe/}{Loupe}
（它展示 JavaScript 的事件循环如何工作）这样的工具，都是有用的教学辅助。
然而，
人们从「自己动手构造可视化」中学到的，
比从「观看别人构造的可视化」中学到的更多~\cite{Stas1998,Ceti2016}，
那么可视化真的有助于学习吗？

为了回答这个问题，
\cite{Cunn2017}重复了一项早先关于学习者在追踪代码执行时做何种草图的研究。
他们发现，完全不做草图与较低的成功率相关，
而在追踪变量值变化时，在变量名旁边写下新值，是最有效的策略。

他们检查过的一个可能的混杂效应是时间：
既然做草图的人解决问题明显花更多时间，
他们是不是仅仅因为想得更久才做得更好？
答案是否定的：
花的时间和得分之间没有相关性。
因此我们的建议是\recommendation{教学习者在调试时追踪变量的值}。

\begin{aside}{流程图}\index{流程图}
  关于可视化，一个常被忽略的发现是：
  \emph{如果两者结构同样良好}，人们理解流程图比理解伪代码更好~\cite{Scan1989}。
  早先那项显示伪代码胜过流程图的研究，用的是结构化伪代码和杂乱流程图；
  把条件拉平之后，
  新手在用图形表示时表现更好。
\end{aside}

\seclbl{我们还能做些什么来帮忙？}{s:pck-help}

\cite{Viha2014}考查了编程课上各种干预措施对通过率的平均提升。
他们指出，有很多理由要对他们的发现打点折扣：
变化前的教学实践很少被清楚说明、
变化的质量没有被评判、
而且只有 8.3\% 的研究报告了负面发现，
所以要么存在正向报告偏差，
要么我们现在教的方式已经是最糟的、任何改变都算改进。
而且和本章讨论的许多其他研究一样，
他们只看大学课堂，
所以他们的发现可能无法推广到其他群体。

记住这些告诫，
他们发现了老师可以做的十件能改善结果的事（\figref{f:pck-interventions}）：

\begin{description}

\item[协作：]
  鼓励学习者在课堂或实验室里协作的活动。

\item[内容调整：]
  部分教学材料被更换或更新。

\item[语境化：]
  课程内容和活动被对齐到一个特定语境，比如游戏或媒体。

\item[CS0：]
  创建一门预备课程，在入门编程课之前修；
  可以只面向某些（例如有风险的）学习者组织。

\item[游戏主题：]
  给课程引入一个以游戏为主题的组成部分。

\item[评分方案：]
  改变评分方案，
  比如提高编程活动的权重、同时降低考试的权重。

\item[小组作业：]
  增加小组作业投入的活动，比如团队式学习和合作学习。

\item[媒体计算：]
  明确声明使用媒体计算的活动（\chapref{s:motivation}）。

\item[同伴支持：]
  以结对、小组、雇用的同伴导师或助学者的形式提供的同伴支持。

\item[其他支持：]
  一个涵盖所有支持活动的总称，
  例如增加教师工时、额外的支持渠道等。

\end{description}

\figimg{figures/interventions-scaled.png}{干预措施的有效性}{f:pck-interventions}{}

这份清单凸显了合作学习的重要性。
\cite{Beck2013}专门在三个学年、由两位不同教师讲授的课程里考察了这一点，
发现总体上以及许多子群体里都有显著收益。
参与合作的人成绩更高，
在期末考试里留空的题目更少，
这表明他们有更高的自我效能感、也更愿意尝试去调试。

\begin{aside}{不写代码的计算}
  写代码并不是教人编程的唯一方式：
  让新手做计算创造力练习，能在好几个层面上提高成绩~\cite{Shel2017}。
  一个典型的练习，是用输入、输出和功能来描述一件日常物品（比如回形针或牙刷）。
  这类教学有时被叫做\grefdex{g:cs-unplugged}{无需计算机的计算机科学}{无需计算机的计算机科学}；
  \hreffoot{https://csunplugged.org/en/}{CS~Unplugged} 网站上就有用于此的课程和练习。
\end{aside}

\seclbl{接下来往哪走？}{s:pck-final}

对想深入下去的人，
\cite{Finc2019}是对计算机教育研究的全面综述，
\cite{Ihan2016}总结了研究最常使用的方法。
我希望有朝一日我们能有像~\cite{Ojos2015}那样的目录，
以及更多像~\cite{Hazz2014,Guzd2015a,Sent2018}这样的教师培训材料，
来帮我们所有人做得更好。

本章报告的研究大多是公共资助的，
却被锁在付费墙之后：
据估计，
我在写这本书时，为了从 \hreffoot{https://en.wikipedia.org/wiki/Sci-Hub}{Sci-Hub} 之类的网站下载论文，违法了 250 次。\index{非法下载（本书作者的）}
我希望不需要这么做的日子早日到来；
如果你是研究者，
请把作品发表在开放获取的渠道上，好让那一天来得更快。

\seclbl{练习}{s:pck-exercises}

\exercise{你的学习者的误解}{小组}{15}

以小组为单位，
重读\secref{s:pck-misunderstand}，列一份你认为你的学习者所持误解的清单。
它们有多具体？
你会怎么检查这份清单有多准确？

\exercise{检查常见错误}{个人}{20}

这些常见错误取自~\cite{Sirk2012}里更长的一份清单：

\begin{description}

\item[赋值方向反了：]
  学习者把左边变量的值赋给了右边变量，
  而不是反过来。

\item[分支判断错了：]
  学习者以为 \texttt{if} 主体里的代码
  即使条件为假也会运行。

\item[执行函数而不是定义它：]
  学习者以为函数在被定义时就被执行。

\end{description}

\noindent
为每一个各写一道练习，用来检查学习者\emph{没有}犯那个错误。

\exercise{乱序代码}{两人一组}{15}

\cite{Chen2017}描述了这样一类练习：学习者重建被弄乱的代码，
弄乱的办法包括删掉注释、
删掉或替换代码行、
移动行序，等等。
这类题目的表现与「让学习者写代码」的评估表现高度相关，
但批改这类题目的工作量更小。
拿你过去出过的一道编程题的解答，
用两种不同的方式把它弄乱，
和同伴交换，
看看你们各自要多久才能正确答出对方的问题。

\exercise{降雨问题}{两人一组}{10}

\cite{Solo1986}提出了降雨问题，
它被用于其后许多编程研究~\cite{Fisl2014,Simo2013,Sepp2015}。
写一个程序，反复读入正整数，直到读到整数 99999。
看到 99999 之后，
程序打印此前读到的所有数字的平均值。

\begin{enumerate}

\item
  用你选择的编程语言解决降雨问题。

\item
  把你的解与同伴的比较。
  你的解做到了哪些他的没做到？反过来呢？

\end{enumerate}

\exercise{变量的角色}{两人一组}{15}

\cite{Kuit2004,Byck2005,Saja2006}提出了一组单变量模式，
我发现它们在教初学者时非常有用：

\begin{description}

\item[固定值：]
  一个在初始化之后就不再取得新的正当值的数据项。

\item[步进器：]
  一个以系统、可预测的方式依次经过一串值的数据项。

\item[遍历器：]
  一个在数据结构中穿行的数据项。

\item[最近持有者：]
  一个持有在依次经过一串值的过程中最近遇到的值的数据项。

\item[最想要持有者：]
  一个持有到目前为最好或最合适值的数据项。

\item[收集器：]
  一个累积各个值之效果的数据项。

\item[跟随者：]
  一个总是从另一个数据项的旧值取得新值的数据项。

\item[单向标志：]
  一个双值数据项，一旦值被改变，就不能再回到初始值。

\item[临时变量：]
  一个只极短暂持有某个值的数据项。

\item[组织器：]
  一个存储可重新排列元素的数据结构。

\item[容器：]
  一个存储可增删元素的数据结构。

\end{description}

选一个 5—15 行的程序，用这些类别给它的变量分类。
把你的分类与同伴的比较。
在你们不一致的地方，
你们理解对方的看法了吗？

\exercise{你在教什么？}{个人}{10}

把你教的主题与~\cite{Luxt2017}（\secref{s:pck-now}）里列出的清单作比较。
你覆盖了哪些主题？
哪些\emph{没有}覆盖？
你覆盖了哪些他们清单里没有的额外主题？

\exercise{有益的活动}{个人}{10}

看看\cite{Viha2014}（\secref{s:pck-help}）列出的干预措施。
其中哪些你在自己的课堂上已经在做？
哪些你能轻易加上？
哪些无关？

\exercise{误解与挑战}{小组}{15}

\hreffoot{http://www.pd4cs.org/}{计算机科学原理教学专业发展}网站
包含\hreffoot{http://www.pd4cs.org/mc-index/}{一份详细的学习者误解与练习清单}。
以小组为单位，
选一个部分（比如数据结构或函数），过一遍他们的清单。
这些误解里，哪些你记得自己当学习者时也有过？
哪些你现在还有？
哪些你在自己的学习者身上见过？

\exercise{你最关心什么？}{全班}{15}

\cite{Denn2019}请人们提出并给各种计算机教育研究问题打分，
发现研究者最关心的问题，
和非研究者最关心的问题之间没有任何重叠。
研究者最喜欢的是：

\begin{enumerate}

\item
  哪些基础编程概念对学生来说最具挑战性？

\item
  在入门编程课上面临广泛的先备经验差异时，哪些教学策略最有效？

\item
  是什么影响了学生从简单编程示例中进行概括的能力？

\item
  对儿童教计算时，哪些教学实践最有效？

\item
  编程课上的学生对哪类问题最投入？

\item
  把编程教给各类群体，最有效的方式是什么？

\item
  把计算教育推广到覆盖全体学生，最有效的方式是什么？

\end{enumerate}

\noindent
而对非研究者来说最重要的问题是：

\begin{enumerate}

\item
  为了改善学习，什么时候、以何种方式给学生的代码提供反馈最好？

\item
  教学生计算机科学时，哪类编程练习最有效？

\item
  对学习计算的学生来说，项目式学习、讲授和主动学习各自的相对优点是什么？

\item
  在编程课上给学生提供反馈最有效的方式是什么？

\item
  人们在编程时把问题拆解成更小任务的过程中，觉得最难的是什么？

\item
  在入门计算课上，学生需要理解的关键概念是什么？

\item
  在非计算学科的学生中培养计算能力，最有效的方式是什么？

\item
  教授基础计算概念和技能的最佳顺序是什么？

\end{enumerate}

让班上每个人独立地，
从合并后的清单里挑出他们最关心的八个问题各投一票，
然后计算每个问题的平均分。
在你的班里哪些最受欢迎？
最受欢迎的问题属于哪一组？
